%!TeX encoding = UTF-8
\documentclass[12pt,a4paper]{article}
\usepackage[utf8]{inputenc}
\usepackage[T1]{fontenc}
\usepackage[spanish,english]{babel}
\usepackage{graphicx}
\usepackage[margin=2.5cm]{geometry}
\usepackage{hyperref}
\usepackage{booktabs}
\usepackage{amsmath}
\usepackage{longtable}

\title{Segmentación Semántica de Imágenes Histológicas de Placenta de Ratón Usando YOLO}
\author{Agustín Páez\\Director: [NOMBRE DEL DIRECTOR]\\Universidad Nacional de San Juan}
\date{\today}

\begin{document}
\maketitle

\begin{abstract}
Este trabajo presenta un sistema de segmentación semántica para imágenes histológicas de placenta de ratón (Whole Slide Images - WSI) utilizando YOLO. El dataset fue preparado manualmente a partir de 50 WSIs anotadas en QuPath, representando tres clases principales: Decidua, Labyrinth y Junctional Zone. Se experimentó con múltiples configuraciones de tamaño de parche (256, 512, 768, 1280, 2048 píxeles) y diferentes niveles de solapamiento (50\%, 150, 256). El mejor modelo (YOLOv11x-seg) alcanzó un mIoU global de 0.78 y Dice de 0.83 en validación.
\end{abstract}

\section{Introducción}

La placenta es un órgano crítico en el desarrollo embrionario de los mamíferos. En modelos ratón, presenta tres regiones principales:
\begin{itemize}
    \item \textbf{Decidua}: Capa exterior que contacta con el útero materno
    \item \textbf{Labyrinth}: Región central de intercambio sanguíneo
    \item \textbf{Junctional Zone}: Zona de transición entre decidua y labyrinth
\end{itemize}

El análisis automático de estas estructuras es relevante para estudios de biología del desarrollo, toxicología reproductiva e investigación de trastornos del embarazo. Sin embargo, la segmentación manual de WSIs es un proceso extremadamente consumidor de tiempo y propenso a errores.

\section{Trabajo Relacionado}

\subsection{Métodos Clásicos}
\begin{itemize}
    \item Thresholding y clustering
    \item Region growing y watershed
    \item Active contours, level sets
\end{itemize}

\subsection{Aprendizaje Profundo para Segmentación}
\begin{itemize}
    \item U-Net: estándar en segmentación médica
    \item DeepLab, SegNet
    \item Transformers: ViT, Swin-UNETR
    \item YOLOv8/YOLOv11-Seg: segmentación en tiempo real
\end{itemize}

\subsection{Datasets Relacionados}
\begin{itemize}
    \item FeTA: 90 imágenes MRI fetal con 7 clases
    \item PlaNet-S: 42 imágenes histológicas de placenta
    \item Mouse embryo: 200 imágenes de embrión ratón
    \item Intestinal organoids: 840 imágenes de organoides ratón
\end{itemize}

\section{Metodología}

\subsection{Dataset}

Se utilizaron 50 Whole Slide Images (WSIs) de placenta de ratón en formato TIF, con dimensiones típicas de ~8000×6000 píxeles.

\begin{table}[ht]
\centering
\caption{Distribución de WSIs utilizadas}
\label{tab:wsi-dist}
\begin{tabular}{@{}lcc@{}}
\toprule
Conjunto & WSIs & Parches generados \\ \midrule
Entrenamiento & 35 & 26,038 \\
Validación & 15 & 11,340 \\ \midrule
\textbf{Total} & \textbf{50} & \textbf{37,378} \\ \bottomrule
\end{tabular}
\end{table}

\subsection{Clases de Anotación}

\begin{table}[ht]
\centering
\caption{Clases del dataset}
\label{tab:classes}
\begin{tabular}{@{}ccc@{}}
\toprule
ID & Clase & Descripción \\ \midrule
0 & Decidua & Capa exterior de la placenta \\
1 & Labyrinth & Región central de intercambio \\
2 & Junctional Zone & Zona de transición \\
3 & no\_labeled & Áreas no anotadas (ignoradas) \\
\bottomrule
\end{tabular}
\end{table}

\subsection{Preprocesamiento y Generación de Parches}

Se generaron múltiples datasets con diferentes configuraciones:

\begin{longtable}{|l|c|c|c|}
\hline
\textbf{Dataset} & \textbf{Tamaño parche} & \textbf{Solapamiento} & \textbf{Clases} \\ \hline
\endfirsthead
\hline
\textbf{Dataset} & \textbf{Tamaño parche} & \textbf{Solapamiento} & \textbf{Clases} \\ \hline
\endhead
\hline
\endfoot
\hline
\endlastfoot
YOLO-256x256-50ovlp & 256 & 50\% & 3 \\
YOLO-256x256-50ovlp\_V2 & 256 & 50\% & 5 (con background) \\
YOLO-256x256-50ovlp\_V3 & 256 & 50\% & 5 \\
YOLO-256x256-50ovlp\_V4 & 256 & 50\% & 4 (sin no\_labeled) \\
YOLO-512x512-150ovlp & 512 & 150px & 3 \\
YOLO-768x768-150ovlp & 768 & 150px & 3 \\
YOLO-1280x1280-256ovlp & 1280 & 256px & 3 \\
YOLO-1280x1280-256ovlp\_w\_bgs & 1280 & 256px & 3 \\
YOLO-2048x2048-512ovlp & 2048 & 512px & 3 \\
\end{longtable}

\subsection{Modelos Entrenados}

Se experimentó con múltiples variantes de YOLO:

\begin{longtable}{|l|c|c|c|}
\hline
\textbf{Modelo} & \textbf{Parámetros} & \textbf{GFLOPs} & \textbf{Dataset} \\ \hline
\endfirsthead
\hline
\textbf{Modelo} & \textbf{Parámetros} & \textbf{GFLOPs} & \textbf{Dataset} \\ \hline
\endhead
\hline
\endfoot
\hline
\endlastfoot
YOLOv8x-seg & 70.5M & 336.7 & YOLO-256x256-50ovlp\_V4 \\
YOLOv8l-seg & 43.7M & 165.4 & YOLO-1280x1280-256ovlp \\
YOLOv8l-seg & 43.7M & 165.4 & YOLO-1280x1280-256ovlp\_w\_bgs \\
YOLOv11x-seg & 62.0M & 295.9 & YOLO-1280x1280-256ovlp\_w\_bgs \\
YOLOv8x-seg & 70.5M & 336.7 & YOLO-512x512-150ovlp \\
YOLOv8x-seg & 70.5M & 336.7 & YOLO-768x768-150ovlp \\
YOLOv8x-seg & 70.5M & 336.7 & YOLO-2048x2048-512ovlp \\
\end{longtable}

\subsection{Configuración de Entrenamiento}

\begin{itemize}
    \item Modelo base: YOLOv8x-seg.pt / YOLOv11x-seg.pt
    \item Epochs: 100
    \item Optimizer: AdamW
    \item Learning rate: 0.001
    \item Weight decay: 0.0005
    \item Momentum: 0.937
    \item Batch size: 128
    \item Patience: 15
    \item Augmentations: Mosaic 1.0, flip horizontal 0.5, rotación ±10°
\end{itemize}

\section{Resultados}

\subsection{Métricas de Evaluación (Modelo YOLOv11x-seg - Mejor resultado)}

\begin{table}[ht]
\centering
\caption{Resultados globales - YOLOv11x-seg}
\label{tab:results-global}
\begin{tabular}{@{}lcccc@{}}
\toprule
Métrica & Valor & std \\ \midrule
mIoU & 0.780 & 0.277 \\
Dice & 0.835 & 0.272 \\
Precision & 0.833 & 0.289 \\
Recall & 0.866 & 0.261 \\
Pixel Accuracy & 0.906 & 0.150 \\
F1 & 0.835 & 0.272 \\ \bottomrule
\end{tabular}
\end{table}

\begin{table}[ht]
\centering
\caption{Resultados por clase - YOLOv11x-seg}
\label{tab:results-per-class}
\begin{tabular}{@{}lcccc@{}}
\toprule
Clase & Box mAP50 & Seg mAP50 & Seg mAP & AP \\ \midrule
Decidua & 0.694 & 0.618 & 0.352 & 0.518 \\
Labyrinth & 0.846 & 0.846 & 0.698 & 0.733 \\
Junctional Zone & 0.646 & 0.635 & 0.393 & 0.476 \\ \bottomrule
\end{tabular}
\end{table}

\subsection{Ultralytics Metrics}

\begin{table}[ht]
\centering
\caption{Métricas de Ultralytics}
\label{tab:ultralytics-metrics}
\begin{tabular}{@{}lcccc@{}}
\toprule
Métrica & Global & Decidua & Labyrinth & JZ \\ \midrule
Box mAP50 & 0.729 & - & - & - \\
Box mAP & 0.576 & - & - & - \\
Seg mAP50 & 0.700 & 0.618 & 0.846 & 0.635 \\
Seg mAP & 0.481 & 0.352 & 0.698 & 0.393 \\ \bottomrule
\end{tabular}
\end{table}

\section{Análisis de Resultados}

\subsection{Por clase}
\begin{itemize}
    \item \textbf{Labyrinth}: Mejor rendimiento (mAP50 = 0.846) - textura más distintiva
    \item \textbf{Decidua}: Rendimiento medio (mAP50 = 0.618) - mayor variabilidad
    \item \textbf{Junctional Zone}: Menor rendimiento (mAP50 = 0.635) - zona de transición ambiguo
\end{itemize}

\subsection{Por tamaño de imagen}
\begin{itemize}
    \item Imágenes mayores (1280+) capturan mejor las estructuras grandes
    \item Imágenes menores (256) son más rápidas pero pierden detalle
\end{itemize}

\section{Conclusiones}

\begin{itemize}
    \item YOLO es efectivo para segmentación de placenta de ratón
    \item Labyrinth es la clase con mejor detección debido a su textura distintiva
    \item Junctional Zone es la clase más difícil por ser una zona de transición
    \item El solapamiento de parches mejora la continuidad de las máscaras
\end{itemize}

\textbf{Palabras clave}: Segmentación semántica, placenta de ratón, YOLO, imágenes histológicas, Whole Slide Images, deep learning

\section{Referencias}

\begin{thebibliography}{99}

\bibitem{yolo-fetal-2026}
YOLOv8-Seg for Multi-Tissue Segmentation of Fetal Brain MRI: A FeTA Benchmark and Comparative Study with U-Net Variants.
\textit{Journal of Imaging Informatics in Medicine}, 2026.

bibitem{planets}
S. Yamamoto, S. Saito.
PlaNet-S: Automatic Semantic Segmentation of Placenta.
\textit{arXiv:2312.11580}, 2023.

bibitem{embryo}
S.M. Rolfe, A.M. Maga.
Deep Learning Enabled Multi-Organ Segmentation of Mouse Embryos.
\textit{bioRxiv}, 2022.

bibitem{organoids}
T. Mäkelä, E. Domenech-Moreno.
Annotated mouse intestinal organoid dataset (YOLO format).
\textit{Zenodo}, 2023.

bibitem{qupath}
Bankhead, P., et al.
QuPath: Open source software for digital pathology.
\textit{Scientific Reports}, 2017.

bibitem{yolo-ultralytics}
Glenn Jocher.
Ultralytics YOLOv8/YOLOv11.
\textit{GitHub}, 2023-2024.

bibitem{unet}
Ronneberger, O., et al.
U-Net: Convolutional Networks for Biomedical Image Segmentation.
\textit{MICCAI}, 2015.

\end{thebibliography}

\end{document}
